\section{DGF source correspondence}
\label{app:dgf}

The source is \citet{canale2026}, a 124-page unpublished practitioner framework. The paper uses it only
for the facts reproduced in Section~\ref{sec:routes-worked} and located in Table~\ref{tab:sources}: the
three initial situations, the opinions recorded on them, the bounded acquisition exception, the route
types, and the ordinal risk scale. Nothing else in the framework is needed to follow or to check the
argument. Page numbers refer to the printed source. The gate types and example tasks of Table~\ref{tab:gates}, the
life-cycle positioning of Table~\ref{tab:lifecycle}, and the illustrative reviews of
Table~\ref{tab:rolecases} summarize the author's governance practice; they are not taken from the pages
listed below. The nine populations, the 140-person allocation, route volumes,
labor-displacement fractions, support floors, calendar scenarios, and proposed machine operations are
this paper's analytical additions. They are not measured staffing or performance data from the
framework.

\begin{table}[!htb]
\centering\small
\caption{Framework locations used by the argument and the three worked routes.}
\label{tab:sources}
\begin{tabularx}{\linewidth}{@{}P{3.4cm}L@{}}
\toprule
Printed pages & Material used\\
\midrule
14, 16 & Dispositions and responsible functions referred to in the three cases.\\
17--20 & The three route types (new platform, acquisition, new project) abstracted here as Buy,
Integrate, and Build.\\
21--23, 28, 32--36 & Required evidence, dated obligations, procurement exchanges, shared IT/Security dossier.\\
38, 90, 98, 100--101 & Business risk ownership, operational handover, ordinal risk scale, exception process.\\
103--105 & CRM red opinion, bounded acquisition exception, support-agent red opinion.\\
\bottomrule
\end{tabularx}
\end{table}

Shared work, cross-gate negotiation, and renewed review require additional visit records in an
empirical implementation.
